%% ========================================================================
\section{Discussion}
\label{sec:discussion}
%% ========================================================================

\subsection{Significance of the Fully Autonomous Workflow}

The patient journey presented in this paper represents a milestone in autonomous AI-driven clinical trial illustration. The 13-commit, 72-minute workflow executed by Claude Code Opus 4.6 produced approximately 8,500 lines of Python code, 30 progress diagrams, 6 deliverable diagrams, 6 regulatory tables, 4 guidance documents, an FDA cost-savings analysis module, and 262 tests - all without user intervention between commits. This level of autonomous output generation in a regulated healthcare domain has not been previously demonstrated.

The significance extends beyond the volume of generated artifacts. Each artifact maintains clinical coherence with every other artifact in the system. The \texttt{patient\_state.py} data model provides the shared vocabulary that ensures patient demographics remain consistent across all 10 stage scripts. The \texttt{master\_journey.py} orchestrator's \texttt{STAGE\_REGULATORY\_MAP} and \texttt{STAGE\_MODULE\_MAP} dictionaries enforce systematic regulatory coverage and module activation across stages. The 30 progress diagrams each show the cumulative state of the patient journey at that point, creating a verifiable chain of clinical progress that any reviewer can inspect stage-by-stage.

The explainability of the autonomous workflow is particularly noteworthy. Unlike a monolithic code generation that produces a single large output, the commit-by-commit approach enabled real-time monitoring of the AI's work product. The text-based diagrams played a critical role in this monitorability: because they are rendered as plain ASCII art, they are immediately viewable in any GitHub branch interface without requiring code execution, specialized rendering tools, or local development environment setup. A clinical reviewer could observe the Stage 3 clinical perspective diagram showing digital twin initialization with four growth models and V\&V 40 validation status, and assess whether the generated content was clinically relevant, without needing to read any Python code.

\subsection{Connection Between Regulatory Adaptations and Patient Journey}

The three regulatory adaptations previously conducted by the author - 21 CFR Part 312 Subpart J (94 pages, 42 references), 21 CFR Part 50 Subpart C (37 pages, 19 references), and ICH E6(R3) - were essential prerequisites for the patient journey to have clinical substance. Without these adapted regulations establishing how Physical AI systems should be governed in oncology clinical trials, the patient journey would have been a technical exercise without regulatory grounding.

The connection is explicit in the codebase. The \texttt{STAGE\_REGULATORY\_MAP} in \texttt{master\_journey.py} maps each stage to specific regulatory sections from all three frameworks. For example, Stage 5 (Robotic Surgery) references \S312.402 for surgical eligibility assessment and ICH E6(R3) \S2.8 for quality management of the robotic procedure. Stage 7 (Adjuvant Therapy) references \S312.404 for treatment safeguards and monitoring, ICH \S2.10 for safety reporting of the three adverse events, and ICH \S2.11 for protocol amendments related to dose modifications. Stage 10 (Data Lock and Archival) references \S312.403 for protocol submission, ICH \S2.9 for data integrity verification at HARD\_LOCK, and ICH \S2.12 for essential document archival with 15-year retention.

This regulatory mapping demonstrates that the adapted Physical AI regulations by Claude Code, based on the original ICH/CFR documents, were necessary for the single patient journey to provide the paper with clinical and regulatory substance. The 84+ regulatory sections addressed across all 10 stages show comprehensive coverage that would not have been possible without the prior regulatory adaptation work.

\subsection{Clinical Implications}

The patient journey illustrates how Physical AI systems can transform each stage of an oncology clinical trial:

\textbf{Pre-screening and enrollment} (Stages 1--2): Digital prescreening with PHI detection, de-identification per \S50.33, FHIR R4 harmonization, and DICOM validation accelerates the traditional screening pipeline. The autonomous workflow generated stage scripts that implement 18 HIPAA identifier checks, ICD-10 to SNOMED mapping, and automated access provisioning with Part 11 signatures.

\textbf{Digital twin and robot qualification} (Stages 3--4): The digital twin engine, initialized with patient-specific tumor models and achieving 94.2\% fidelity, provides treatment response simulation and risk trajectory prediction throughout the journey. Robot qualification via USL scoring (da Vinci Xi: 87.5, Franka Emika: 88.75) ensures that deployed systems meet minimum thresholds across autonomy, dexterity, safety, and interoperability dimensions before first patient contact.

\textbf{Surgery and recovery} (Stages 5--6): The 168-minute robotic lobectomy with 85 mL estimated blood loss, negative margins, and 18 lymph nodes sampled demonstrates the precision capabilities of Physical AI surgical systems. The single safety event (12.1N force spike, resolved in 200ms) shows the real-time monitoring capabilities operating within the safety architecture defined in \texttt{deliverables/diagrams/diagram\_05\_safety\_architecture.txt}.

\textbf{Immunotherapy and federation} (Stages 7--8): The 35-cycle pembrolizumab treatment achieving CR, combined with 70 federated learning rounds across 12 sites using SMPC with differential privacy ($\epsilon = 1.0$), illustrates the integration of active treatment with multi-site collaborative analytics. The federated hazard ratio of 0.62 (95\% CrI: 0.44--0.87) and DSMB recommendation to CONTINUE demonstrate the real-time safety governance operating throughout treatment.

\textbf{Surveillance and closeout} (Stages 9--10): The progressive risk reduction from 35\% to 3\% over four federated model updates, 36-month event-free survival, and comprehensive closeout with HARD\_LOCK, CSR generation, and GCP COMPLIANT audit across all 10 domains show how the journey completes with full regulatory compliance and data integrity.

\subsection{Role of the Diagram Evolution}

The 30-diagram progression across 10 stages demonstrates both clinical and Physical AI proficiency. As the patient advances through each stage, the diagrams evolve to show increasing complexity:

\begin{itemize}[leftmargin=*]
\item \textbf{Stage 1 diagrams}: Show 6 of 51 modules activated, 7 regulatory sections addressed, and MCP conformance at CLINICAL\_READ level.
\item \textbf{Stage 5 diagrams}: Show 38 of 51 modules activated, accumulated surgical metrics (168 min, 85 mL EBL, negative margins, 18 LN), and MCP conformance elevated to ROBOT\_PROCEDURE level.
\item \textbf{Stage 10 diagrams}: Show 51 of 51 modules activated, all 84+ regulatory sections verified, HARD\_LOCK applied, and final patient status COMPLETED with CR, 36-month EFS, and 0 device-related adverse events.
\end{itemize}

This evolution demonstrates the explainability of Claude Code to humans. A reviewer examining the diagrams in sequence can quickly assess whether the autonomous workflow produced clinically coherent progressions. The consistent format across all 30 diagrams - with standardized sections for milestones, cumulative progress, status indicators, and key metrics - enables rapid pattern recognition and anomaly detection, contributing to the credibility of the autonomous process.

\subsection{Safety Architecture and Monitoring}

The monitoring architecture, detailed in \texttt{deliverables/diagrams/diagram\_05\_safety\_architecture.txt}, implements a three-pillar approach: real-time monitoring (vital signs, digital twin state, robot telemetry), adverse event reporting (detection, classification, documentation, expedited reports), and oversight bodies (DSMB, IRB/IEC, FDA). The three-tier alert threshold system (YELLOW: warning deviation, ORANGE: action required, RED: immediate intervention/e-stop) provides graduated escalation appropriate to event severity.

For PAT-2026-0042, the safety architecture managed all three adverse events appropriately: the Grade 2 atrial fibrillation triggered an ORANGE alert leading to dose hold and cardiology consultation, while the Grade 1 hypothyroidism and rash triggered YELLOW alerts managed with standard clinical interventions. The absence of any RED alerts or device-related adverse events throughout the 35-cycle treatment and 36-month follow-up demonstrates the robustness of the Physical AI safety monitoring integration.

\subsection{Comparison with Traditional Trial Conduct}

The Physical AI-augmented trial conduct illustrated in this patient journey differs from traditional approaches in several ways. The \texttt{fda\_cost\_savings\_analysis.py} module quantifies these differences across four dimensions:

\begin{enumerate}[leftmargin=*]
\item \textbf{Cost}: 30--50\% total reduction (\$390M--\$650M), with digital twin simulation contributing the largest savings (15--25\%) through virtual patient cohort optimization, in-silico dose-response modeling, and predictive toxicity simulations.

\item \textbf{Time}: 18--30 months total acceleration, comprising 12--18 months from enrollment acceleration through digital prescreening and AI-driven patient matching, plus 6--12 months from faster safety signal detection via continuous robotic monitoring and federated cross-site analysis.

\item \textbf{Quality}: 30\% reduction in protocol deviations, 25\% faster adverse event detection, and 15\% improvement in data completeness, driven by automated procedure guidance, continuous sensor monitoring, and digital twin data capture.

\item \textbf{Per-patient}: Average savings of \$11,000 per patient (24.4\% reduction from \$45,000 to \$34,000), stemming from reduced screen failures, fewer monitoring visits, and decreased protocol deviation remediation costs.
\end{enumerate}

These projections are subject to variability based on trial design, disease indication, Physical AI system maturity, and regulatory requirements under 21 CFR Part 312 Subpart J, as noted in the cost-savings summary report.

%% ========================================================================
\section{Limitations and Future Work}
\label{sec:limitations}
%% ========================================================================

\subsection{Limitations}

This work has several inherent limitations that should be considered when interpreting the results:

\textbf{Illustrative nature.} The patient journey is a simulated illustration, not a report from an actual clinical trial. PAT-2026-0042 is a modeled patient with representative but synthetic clinical characteristics. No real patient data was used, and no actual robotic procedures, digital twin computations, or federated learning rounds were executed. The clinical outcomes (CR, HR 0.62, 36-month EFS) represent illustrative projections, not empirical results. This publication is based on 3 Physical AI regulatory adaptations conducted by the author and should not be considered a new or approved regulation.

\textbf{Single-patient scope.} The journey follows a single patient through 10 stages. Real oncology trials involve hundreds to thousands of patients with diverse clinical profiles, treatment responses, and adverse event patterns. Extrapolating from a single illustrative journey to trial-level conclusions requires validation with multi-patient simulations and, ultimately, real clinical data.

\textbf{Autonomous generation constraints.} While the 72-minute, 13-commit autonomous workflow demonstrated the capability of Claude Code Opus 4.6 to produce coherent clinical trial infrastructure, the generated code has not been clinically validated. The Python modules implement simulated clinical workflows with placeholder computations rather than validated medical algorithms. The cost-savings projections, while based on published literature baselines, represent modeled estimates subject to significant variability.

\textbf{Regulatory adaptation scope.} The three adapted regulatory frameworks (21 CFR Part 312 Subpart J, 21 CFR Part 50 Subpart C, ICH E6(R3)) were adapted by the author for Physical AI systems but have not undergone formal regulatory review. The adaptations provide a framework for illustrating how existing regulations might be extended to govern Physical AI clinical trials, but they do not carry regulatory authority.

\textbf{Technology readiness.} The Physical AI ecosystem described (da Vinci Xi, Franka Emika, digital twin engine, federated learning network) represents a composite of commercially available and conceptual capabilities. Full integration of these systems at the level described in the patient journey will also require additional engineering, validation, and regulatory clearance.

\subsection{Future Work}

Future work will continue with fully autonomous workflows in Physical AI oncology trials along several dimensions:

\begin{enumerate}[leftmargin=*]
\item \textbf{Multi-patient journey expansion.} Extend the autonomous workflow to generate multi-patient cohort journeys with diverse clinical profiles, treatment arms, and outcome trajectories. This would demonstrate the scalability of the autonomous approach from single-patient illustration to trial-level simulation.

\item \textbf{Enhanced autonomous workflow capabilities.} Increase the sophistication of autonomous code generation to include executable digital twin computations, validated pharmacokinetic/pharmacodynamic models, and functional federated learning implementations with real model training across synthetic datasets.

\item \textbf{Clinical validation partnerships.} Engage with NCI-designated cancer centers and robotic surgery programs to validate the patient journey pipeline against real clinical workflows, identifying gaps between the illustrated pipeline and actual clinical practice.

\item \textbf{Regulatory engagement.} Submit the adapted regulatory frameworks and patient journey illustration to FDA for informal feedback, potentially through the Pre-Submission (Q-Sub) process, to assess alignment between the adapted regulations and the agency's evolving perspective on Physical AI in clinical trials.

\item \textbf{Extended cost-savings validation.} Refine the cost-savings projections in \texttt{fda\_cost\_savings\_analysis.py} with real-world cost data from trials incorporating robotic assistance, digital twin platforms, or federated learning infrastructure.

\item \textbf{International regulatory expansion.} Extend the regulatory mapping in \texttt{master\_journey.py} to include EU MDR, UK MHRA, and Health Canada frameworks, enabling multi-jurisdictional patient journey illustrations.

\item \textbf{Real-time diagram generation.} Develop capabilities for generating progress diagrams in real time during actual real-life trial execution, extending the monitoring framework from retrospective illustration to prospective clinical trial oversight.
\end{enumerate}

%% ========================================================================
\section{Conclusions}
\label{sec:conclusions}
%% ========================================================================

This paper presents the first fully autonomous single-patient journey through a regulated Physical AI oncology clinical trial illustration, demonstrating that advanced AI systems can produce clinically coherent, regulatory-compliant trial infrastructure without human intervention. The journey of PAT-2026-0042 through 10 stages - from pre-screening through trial closeout  - was orchestrated by Claude Code Opus 4.6 in a single pull request of 13 commits completed in 72 minutes, producing 12 Python modules (approximately 8,500 lines), 30 progress diagrams, 6 deliverable diagrams, 6 regulatory tables, 4 guidance documents, an FDA cost-savings analysis, and 262 tests.

The clinical outcomes illustrated for the patient - complete response to treatment, R0 surgical resection via da Vinci Xi robotic lobectomy, 35 completed pembrolizumab immunotherapy cycles, 36-month event-free survival, and recurrence risk reduction from 35\% to 3\% - demonstrate the comprehensive scope of the illustrated journey. The three adapted regulatory frameworks (21 CFR Part 312 Subpart J, 21 CFR Part 50 Subpart C, ICH E6(R3)) provided the regulatory grounding necessary for clinical substance, with 84+ regulatory sections systematically addressed across all 10 stages through the \texttt{STAGE\_REGULATORY\_MAP} in \texttt{master\_journey.py}.

The FDA cost-savings analysis projects 30--50\% total cost reduction (\$390M--\$650M savings) for Physical AI-augmented oncology trials, with 18 - 30 months of timeline acceleration and quality improvements of 15 - 30\% across protocol compliance, adverse event detection, and data completeness. The USL-scored robotic systems (da Vinci Xi: 87.5, Franka Emika: 88.75), digital twin engine (94.2\% fidelity), and federated learning network (70 rounds across 12 sites with differential privacy) illustrate the Physical AI ecosystem that could enable these improvements.

The autonomous workflow's credibility was enhanced by two mechanisms that distinguish it from traditional AI-generated outputs: the commit-by-commit structure enabling real-time inspection of each stage's artifacts, and the text-based diagram system providing human-readable visual summaries viewable directly in the GitHub branch interface without code execution. These mechanisms demonstrate that fully autonomous AI workflows can be made explainable and monitorable, a requirement for any application in regulated clinical domains.

The patient journey represents a convergence of Physical AI technology, adapted clinical trial regulations, and autonomous AI code generation capabilities. As these three domains continue to advance, the framework established in this work - Python-based stage orchestration, regulatory mapping, USL scoring, cost-savings modeling, and multi-perspective progress tracking - provides a foundation for increasingly sophisticated autonomous clinical trial illustrations that can inform the design and conduct of real Physical AI oncology trials. The complete codebase is publicly available in the \texttt{physical-ai-oncology-trials} repository at version 2.7.0.

%% ========================================================================
%% REFERENCES
%% ========================================================================

\begin{thebibliography}{20}

\bibitem{kawchak2026physicalai}
K.~Kawchak, ``Physical AI Oncology Trials,'' GitHub repository, 2026.
\url{https://github.com/kevinkawchak/physical-ai-oncology-trials}. Patient Journey: \url{https://github.com/kevinkawchak/physical-ai-oncology-trials/tree/main/patient-journey}.

\bibitem{kawchak2026cfr312}
K.~Kawchak, ``Adaptation of 21 CFR Part 312 for Physical AI Oncology Clinical Trials,'' Zenodo, 2026.
\url{https://doi.org/10.5281/zenodo.19057628}

\bibitem{kawchak2026cfr50}
K.~Kawchak, ``Adaptation of 21 CFR Part 50 for Physical AI Human Subject Protections,'' Zenodo, 2026.
\url{https://doi.org/10.5281/zenodo.19040707}

\bibitem{kawchak2026iche6r3}
K.~Kawchak, ``Adaptation of ICH E6(R3) for Physical AI Good Clinical Practice,'' Zenodo, 2026.
\url{https://doi.org/10.5281/zenodo.18973368}

\bibitem{kawchak2026usl}
K.~Kawchak, ``Unification Standard Level for Physical AI Oncology Trials,'' Zenodo, 2026.
\url{https://doi.org/10.5281/zenodo.18778220}

\bibitem{kawchak2026national}
K.~Kawchak, ``National MCP Servers for Physical AI Oncology Clinical Trial Systems,'' Zenodo, 2026.
\url{https://doi.org/10.5281/zenodo.18916731}

\bibitem{cfr312}
U.S. Food and Drug Administration, ``21 CFR Part 312 - Investigational New Drug Application,'' Code of Federal Regulations, Title 21, 2024.

\bibitem{cfr50}
U.S. Food and Drug Administration, ``21 CFR Part 50 - Protection of Human Subjects,'' Code of Federal Regulations, Title 21, 2024.

\bibitem{cfr11}
U.S. Food and Drug Administration, ``21 CFR Part 11 - Electronic Records; Electronic Signatures,'' Code of Federal Regulations, Title 21, 2003.

\bibitem{ichgcp2023}
International Council for Harmonisation, ``ICH E6(R3) Guideline for Good Clinical Practice,'' ICH Harmonised Guideline, 2023.

\bibitem{fhir2019r4}
HL7 International, ``Fast Healthcare Interoperability Resources (FHIR) R4,'' HL7 Standard, 2019.
\url{https://www.hl7.org/fhir/}

\bibitem{dicom2024standard}
NEMA, ``Digital Imaging and Communications in Medicine (DICOM) Standard,'' National Electrical Manufacturers Association, 2024.

\bibitem{mcmahan2017fedavg}
B.~McMahan, E.~Moore, D.~Ramage, S.~Hampson, and B.~A.~y~Arcas, ``Communication-Efficient Learning of Deep Networks from Decentralized Data,'' in \textit{Proceedings of AISTATS}, 2017.

\bibitem{eisenhauer2009recist}
E.~A.~Eisenhauer et al., ``New response evaluation criteria in solid tumours: Revised RECIST guideline (version 1.1),'' \textit{European Journal of Cancer}, vol.~45, no.~2, pp.~228--247, 2009.

\bibitem{hipaa1996}
U.S. Congress, ``Health Insurance Portability and Accountability Act of 1996 (HIPAA),'' Law 104-191, 1996.

\bibitem{mcp2024protocol}
Anthropic, ``What is the Model Context Protocol (MCP)?'' Open Protocol Specification, 2024.
\url{https://modelcontextprotocol.io/}

\bibitem{tufts2020cost}
Tufts Center for the Study of Drug Development, ``Cost of Clinical Trials by Phase and Therapeutic Area,'' Tufts CSDD Impact Report, 2020.

\bibitem{fda2019adaptive}
U.S. Food and Drug Administration, ``Adaptive Designs for Clinical Trials of Drugs and Biologics: Guidance for Industry,'' FDA Guidance Document, 2019.

\end{thebibliography}

\vspace{-0.15cm}

%% ========================================================================
\section*{Acknowledgments}
%% ========================================================================
\vspace{-0.2cm}
The author would like to acknowledge Anthropic for providing access to Claude Code for project and paper generations; OpenAI for providing access to ChatGPT for detailed summaries, and Google Gemini for additional clarifications.
\vspace{-0.2cm}
%% ========================================================================
\section*{Ethical Disclosures}
%% ========================================================================
\vspace{-0.2cm}
The author of the article declares no competing interests.
\vspace{-0.2cm}
%% ========================================================================
\section*{Rights and Permissions}
%% ========================================================================
\vspace{-0.2cm}
This article is distributed under the terms of the Creative Commons Attribution 4.0 International License (CC BY 4.0), which permits unrestricted use, distribution, and reproduction in any medium, provided the original author(s) and source are properly credited, a link to the Creative Commons license is provided, and any modifications made are indicated. To view a copy of this license, visit \url{https://creativecommons.org/licenses/by/4.0/}.
\vspace{-0.2cm}
%% ========================================================================
\section*{Cite This Article}
%% ========================================================================
\vspace{-0.2cm}
Kawchak K. A Cancer Patient's Journey Through a Regulated Physical AI Oncology Trial. Zenodo. 2026; \href{https://doi.org/10.5281/zenodo.19119939}{10.5281/zenodo.19119939}.

% \vspace{3cm}

\end{document}
